Baseline fixed-effects estimates show that a 1 percent increase in AI adoption is associated with approximately -0.010 percent higher TFP (p-value below 0.001), controlling for human capital. The estimated AI-TFP relationship remains directionally stable under stronger identification checks, including two-way fixed effects (-0.002, p-value 0.432) and trimmed-sample estimation (-0.010, p-value below 0.001). For macro growth outcomes, the AI coefficient in the GDP growth specification is -0.001 (p-value 0.049), suggesting the productivity channel may emerge before fully translating into contemporaneous GDP growth. Sensitivity checks also support the main findings: the lagged AI specification yields an AI coefficient of -0.002 (p-value 0.446), and inference remains similar under time-clustered standard errors (-0.002, p-value 0.259). The placebo outcome test using log human capital reports 0.000 (p-value 0.776), which should be interpreted as a falsification check rather than a causal estimate. A two-way fixed-effects specification with Driscoll-Kraay standard errors, which is robust to broad cross-sectional dependence, also preserves a positive AI coefficient (-0.002, p-value 0.383).